\section{Execution and Reproducibility}\label{sec:execution}
Each run is a fresh child process (\texttt{python -I}, scrubbed environment, private working directory), wrapped on macOS in \texttt{sandbox-exec} with a deny-network, write-confined profile; resource limits (CPU, file size, open files) are set in the child and a parent watchdog kills the process group (Fig.~\ref{fig:isolation}). The platform was one macOS machine and CPython 3.13.15. Memory limits rely on parent polling; address-space and process-count limits are not enforced on this platform. Reads are not confined, and a child that detaches into a new session could escape the group kill (untested). \texttt{sandbox-exec} is deprecated. Every record carries a flag that the isolation policy is below the specification minimum, which the specification classes as synthetic or benchmark-only grade. Linux and container isolation are not implemented and were not tested. \textbf{Therefore no claim of safe execution of untrusted code, and no cross-platform claim, is made; verified execution is macOS only.}

Reproduction: the unit tests (74 passed, 1 skipped on macOS, as re-run by the reviewer), the K1 spike driver, the baseline driver, the pair-table builder, the statistics script (project-cluster bootstrap $B=10000$, seed 42) and the figure script are listed in \path{research/reproducibility/REPRODUCIBILITY_PROTOCOL.md}. Seeds: bootstrap and sign-flip 42, power simulation 20261117, B3 mutation 1001. Determinism controls: single-threaded numerics, fixed hash seed, at most four workers. Verdicts, levels, primary reasons and witness ids are expected to reproduce exactly; bootstrap intervals to numeric precision; CPU times only approximately. Analysis libraries are not version-pinned in the lock file. This manuscript was not compiled (no LaTeX toolchain), and diagrams are unrendered.
\dotfig{isolation}{Per-run isolation: macOS sandbox child; Linux container NOT IMPLEMENTED}{fig:isolation}
